%%%%

%%%   Самарский государственный технический университет

%%%   Кафедра прикладной математики и информатики

%%%

%%%   Преамбула для статьи в журнал "Вестник Самарского государственного

%%%   технического университета. Серия: «Физико-математические науки»"

%%%   Ориентирована под MikTEx 2.4 for Win32

%%%

%%%   Хотя сам журнал собирается под GNU/Linux на дистрибутиве TeXLive



\documentclass[11pt,a4paper,twoside]{article}



\usepackage[cp1251]{inputenc}

\usepackage[T2A]{fontenc}

\usepackage[english,russian]{babel}

\usepackage{samgtubib}

\usepackage[dvips]{graphicx}

\usepackage[multiple]{footmisc}



\usepackage{samgtu}

\renewcommand{\VestnikYear}{2009} % Год издания Вестника

\renewcommand{\VestnikNumber}{2\,(19)} % Номер Вестника

\renewcommand{\VestnikMonth}{Ноябрь} % Месяц выхода Вестника

\renewcommand{\VestnikData}{01.11.09} % Дата подписания Вестника в печать

\renewcommand{\VestnikUPL}{26,64} % Усл. печ. л.

\renewcommand{\VestnikUIL}{25,73} % Уч.-изд. л.

\renewcommand{\VestnikRegNumber}{479/09} % Регистрационный номер

\renewcommand{\VestnikOrder}{994} % Номер заказа







%

%\usepackage{pst-barcode}

%\usepackage{auto-pst-pdf}

%%

%%\newcommand{\totop}[1]{\raisebox{6pt}[1pt][2pt]{#1}} 

%%\newcommand{\tottop}[1]{\raisebox{12pt}[1pt][2pt]{#1}} 

%%\newcommand{\totttop}[1]{\raisebox{18pt}[1pt][2pt]{#1}} 

%%\newcommand{\tottttop}[1]{\raisebox{24pt}[1pt][2pt]{#1}} 

%%\newcommand{\totttttop}[1]{\raisebox{30pt}[1pt][2pt]{#1}} 

%%\newcommand{\tobot}[1]{\raisebox{-6pt}[1pt][2pt]{#1}} 

%%\newcommand{\tobbot}[1]{\raisebox{-12pt}[1pt][2pt]{#1}} 

%%\newcommand{\tobbbot}[1]{\raisebox{-18pt}[1pt][2pt]{#1}}

%

%\graphicspath{{./00PIC/}}

%

%%\samgtupubl % для генерации pdf, отдаваемого в типографию СамГТУ

%\pdfcrop % обрезка pdf для размещения на math-net.ru

%\draft % На корректуре

%\filllastpage % Последняя страница не заполнена не полностью

%%\briefcomm % Статья является кратким сообщением



\vsgtu{1384}



\FirstPage{634}

\LastPage{649}

%

%

%\begin{document}





\hfill \begin{pspicture}(1in,1in)

\psbarcode{http://dx.doi.org/10.14498/vsgtu1384}{}{qrcode}

\end{pspicture}

\vspace{-1in}



%\Partition{Дифференциальные уравнения}{Differential Equations}



%%%%%%%%%%%%%%%%%%%%%%%%%%%%%%%%%%%%%%%%%%%%%%%%%%%%%%%%%%%%%%%%%%%%%%%%%%%%%%%%%%%%

% Информация о статье и об авторах на русском и английском языках







\udk{517.956.6}%Уравнения смешанного типа

\msc{35M12}% Коды MSC см. http://www.ams.org/msc/





\rutitle{Задача Дирихле для уравнения смешанного типа с~двумя линиями перехода в~прямоугольной области} {Задача Дирихле для уравнения смешанного типа \dots}





\titlehack{Задача Дирихле для уравнения смешанного \\ типа с двумя линиями перехода \\ в прямоугольной области$^*$}



\ruauthor{А.~А.~Гималтдинова}

{Г\,и\,м\,а\,л\,т\,д\,и\,н\,о\,в\,а~А.~А.}



\ruaffil{\rubashgusterl.}



\enaffil{\enbashgusterl.}



\ruauthordetails{1}{\fullauthor{\rupid{11102}{Альфира Авкалевна Гималтдинова}} \regalia{к.ф.-м.н., доц.; \mem{alfiragimaltdinova@mail.ru}}

\position{доцент}\dept{каф. математического анализа}

 \par

\smallskip

\noindent

$^*$Настоящая статья представляет собой расширенный вариант доклада \cite{gimkur:add1-ru}, сделанного авторами на Четвёртой международной конференции «Математическая физика и её приложения» (Россия, Самара, 25 августа -- 1 сентября 2014).}





\enauthordetails{1}{\fullauthor{\enpid{11102}{Alfira A. Gimaltdinova}} \regalia{Cand. Phys. \& Math. Sci.; \mbox{\mem{alfiragimaltdinova@mail.ru}}}\position{Associate Professor}

\dept{Dept. of Mathematical Analysis)}

\par

\smallskip

\noindent

$^*$This paper is an extended version of the paper \cite{gimkur:add1-eng}, presented at the Mathematical Physics and Its Applications 2014 Conference.

}



\rukeywords{уравнение смешанного типа, задача Дирихле, единственность, существование решения, спектральный метод}



\ruabstract{Изучена первая краевая задача для уравнения эллиптико"=гиперболического типа с двумя перпендикулярными внутренними линиями изменения типа и спектральным параметром. Доказаны единственность и существование решения. При доказательстве единственности используется полнота в пространстве $L_2$ системы, биортогонально сопряженной с~системой собственных функций  соответствующей одномерной задачи. При построении решения в виде суммы ряда по биортогональной системе функций возникает проблема малых знаменателей. Получены оценки об отделимости знаменателей от нуля.}



\entitle{The Dirichlet problem for mixed type equation  with two lines  of degeneracy  in a rectangular area}



\engtitlehack{The Dirichlet problem for mixed type equation  with two lines  of degeneracy  in a rectangular area$^*$}

%

\enauthor{A.~A.~Gimaltdinova}{G\,i\,m\,a\,l\,t\,d\,i\,n\,o\,v\,a~A.~A.}





\enabstract{We study the first boundary value problem for the elliptic-hyperbolic type equation with two perpendicular lines of change of type and spectral parameter. We prove the existence and uniqueness of the solution. In the proof of the uniqueness of solution we use the completeness of biorthogonal system in space $ L_2$. When building a solution as the sum of a series there is a problem of small denominators. We obtained estimates of the denominators of the separation from zero.}



\enkeywords{mixed type equation, Dirichlet problem, uniqueness, existence of a solution, spectral method}



\date{19/XII/2014}

\revisiondate{20/II/2015}

\accepteddate{08/IV/2015}



%%%%%%%%%%%%%%%%%%%%%%%%%%%%%%%%%%%%%%%%%%%%%%%%%%%%%%%%%%%%%%%%%%%%%%%%%%%%%%%%%%%%





\makerutitle



%Каждый пункт статьи должен начинаться с команды Название пункта

%после названия точку ставить не нужно. Пункты автоматически нумеруются.

%Если пункт нумеровать не нужно, то следует использовать в команде

% \Section необязательный параметр [n]:

%   \Section[n]{Имя пункта}

% Введение и заключение обычно не нумеруются.



{\bf 1. Постановка задачи.}

Рассмотрим уравнение

\begin{equation}

Lu\equiv \mathop{\rm sgn} x \cdot  u_{xx}+\mathop{\rm sgn} y \cdot  u_{yy}+\lambda u=0,

\label{gim:eq-uravn1}

\end{equation}

где $\lambda\in \mathbb{C},$ в  области 

$$D=\{(x,y)\in \mathbb{R}^2\,|\,-1< x< 1,\,-\alpha< y< \beta\},$$ 

$\alpha,\beta\in \mathbb{R}$; $\alpha,\beta>0.$

Обозначим $D_1=D\cap\{x>0,y>0\},$ $D_2=D\cap\{x>0,y<0\},$

$D_3=D\cap\{x<0, y<0\},$ $D_4=D\cap\{x<0, y>0\}.$



\smallskip



{\small \sc Задача Дирихле (Задача $D$).} 

{\it Найти функцию $u(x,y),$ удовлетворяющую

условиям}

\begin{gather}

u(x,y)\in {\rm C}(\,\overline D\,)\cap {\rm C}^1(D)\cap {\rm C}^2

(D_1\cup D_2\cup D_3\cup D_4),\label{gim:eq-u2}\\

Lu(x,y)\equiv 0,\;\;(x,y)\in D_1\cup D_2\cup D_3\cup D_4,\label{gim:eq-u5}\\

u(x,y)\big|_{x=1}=u(x,y)\big|_{x=-1}=0,\;\;y\in [-\alpha;\beta],\label{gim:eq-u5}\\

u(x,y)\big|_{y=-\alpha}=\psi(x),\qquad u(x,y)\big|_{y=\beta}=\varphi(x),\;\;x\in [-1;1],\label{gim:eq-u5}

\end{gather}

\it где $\varphi$ и $\psi$ "--- заданные достаточно гладкие функции$,$  $\varphi(\pm 1)=\psi(\pm 1)=0.$

\rm 



\smallskip



Краевые задачи для уравнений смешанного типа с одной или несколькими линиями изменения или вырождения типа были объектом изучения многих авторов.



В работе \cite{gim:bib:mon} построена теория задачи Трикоми для уравнений смешанного типа в классической смешанной области, в которой гиперболическая часть состоит из двух подобластей, ограниченных характеристиками уравнения и линиями изменения типа. Там приведен достаточно полный обзор работ, посвященных данному направлению.



В работе \cite{gim:bib:rasias} предложена задача для уравнения с двумя линиями вырождения в смешанной области, состоящей из четырех эллиптических подобластей и четырех гиперболических подобластей, последние из которых ограничены характеристиками данного уравнения и линиями изменения типа.



В \cite[с.~303]{gim:bib:frankl} было показано, что некоторые задачи трансзвуковой газовой динамики сводятся к задаче Дирихле для уравнений смешанного типа. В \cite{gim:bib:bits} доказана некорректность задачи Дирихле для уравнения Лаврентьева. 

В~дальнейшем задача Дирихле для уравнений смешанного типа привлекала внимание многих авторов [\citen{gim:bib:shabat}--\citen{gim:bib:sold2}]. 

В~этих работах  единственность решения задачи Дирихле для уравнений смешанного типа с одной линией вырождения или изменения типа доказана с помощью принципа экстремума или методом интегральных тождеств, а существование "--- методом интегральных уравнений или разделения переменных.



В работах \cite{gim:bib:sab, gim:bib:sab-vag} исследована задача Дирихле для уравнения смешанного типа с одной внутренней линией степенного вырождения и вырождением на границе в прямоугольной области и методами спектрального анализа установлен критерий единственности и решение задачи построено в виде суммы ряда по системе собственных функций.



В данной работе впервые для уравнения \eqref{gim:eq-uravn1} с двумя внутренними перпендикулярными линиями изменения типа со спектральным параметром изучена задача Дирихле в прямоугольной области $D.$ 

Найдены собственные значения соответствующей спектральной задачи. 

Установлен критерий единственности, и решение задачи \eqref{gim:eq-u2}--\eqref{gim:eq-u5} построено в виде суммы ряда по биортогональной системе функций. 

Единственность решения поставленной задачи доказана на основании полноты биортогональной системы в пространстве $L_2[-1, 1].$ 

Ранее такая идея использовалась в работе \cite{gim:bib:iljin} при доказательстве единственности решения начально-граничной задачи для гиперболических уравнений и в работе \cite{gim:bib:sab} для уравнений смешанного типа с одной линией изменения типа. 

При доказательстве существования решения задачи \eqref{gim:eq-u2}--\eqref{gim:eq-u5} аналогично \cite{gim:bib:arnold, gim:bib:lomov2, gim:bib:sab} возникла так называемая <<проблема малых знаменателей>>, которая создает трудности при обосновании сходимости построенного ряда в классе функций \eqref{gim:eq-u2}. 

При определенных ограничениях на параметры $\alpha,$ $\beta$ доказаны леммы об отделимости малых знаменателей от нуля.



\smallskip



\pagebreak



{\bf 2. Спектральная задача. Построение биортогональной системы.}



\smallskip



{\small \sc Задача $ D_\lambda$.} \it Найти значения $\lambda$ и соответствующие им функции $u(x,y),$ удовлетворяющие задаче \eqref{gim:eq-u2}--\eqref{gim:eq-u5}$,$ где $\varphi(x)=\psi(x)\equiv 0.$ \rm



\smallskip



После разделения переменных $u(x,y)=X(x)Y(y)$ получим два обыкновенных дифференциальных уравнения:

\begin{gather}

\mathop{\rm sgn}x\cdot X''+d\cdot X=0,\quad x\in (-1, 0)\cup (0, 1),\label{gim:eq-uravn2}\\

\mathop{\rm sgn}y\cdot Y''+(\lambda-d)Y=0,\quad y\in (-\alpha, 0)\cup(0, \beta),\label{gim:eq-uravn3}

\end{gather}

где $d=\mu^2\in \mathbb{C}$ "--- постоянная разделения, причем в силу условий задачи должны быть выполнены условия

\begin{gather}

X(0-0)=X(0+0),\quad X'(0-0)=X'(0+0),\label{gim:eq-u6}

\\

X(-1)=X(1)=0,\label{gim:eq-uu6}

\\

Y(0-0)=Y(0+0),\quad Y'(0-0)=Y'(0+0),\label{gim:eq-u7}

\\

Y(-\alpha)=Y(\beta)=0.\label{gim:eq-uu7}

\end{gather}



Спектральные задачи \eqref{gim:eq-uravn2}, \eqref{gim:eq-u6}, \eqref{gim:eq-uu6} и  \eqref{gim:eq-uravn3}, 

\eqref{gim:eq-u7}, \eqref{gim:eq-uu7} не являются классическими из-за незнакоопределенности коэффициента при старшей производной. %В работе \cite{pjat}  изучены вопросы о базисности Рисса собственных и присоединенных функций таких задач.



Решениями уравнения \eqref{gim:eq-uravn2}, удовлетворяющими  условиям \eqref{gim:eq-u6}, являются функции

$$

X(x)=\begin{cases}

C_1\cos \mu x+C_2\sin \mu x, & x>0,\\ 

C_1\ch \mu x+C_2\sh \mu x, & x<0,

\end{cases}$$

где $C_1,$ $C_2$ "--- произвольные постоянные, и с учетом условия \eqref{gim:eq-uu6} для $\mu$ получим уравнение

\begin{equation}

\tg \mu=-\th \mu.

\label{gim:eq-uravnmu1}

\end{equation}



\smallskip



\hypertarget{gim:l1}{}



{\small \sc Лемма 1.} {\it Уравнение

\begin{equation}

\tg (az)=-\th (bz), \quad a, b\in {\mathbb R}, \quad a, b> 0,\label{gim:eq-gim:ur-tang}

\end{equation}

имеет счетное множество корней$,$ состоящее из нуля$,$ простых попарно противоположных действительных и попарно противоположных чисто мнимых корней$,$ для которых справедливы следующие асимптотические представления}\/:

$$z_k^{(1),(2)}=\pm\Bigl(-\frac{\pi}{4a}+\frac{\pi}{a}k+O\big(e^{-2\pi kb/a}\big)\Bigr),$$

$$z_k^{(3),(4)}=\pm i \Bigl(-\frac{\pi}{4b}+\frac{\pi}{b}k+O\big(e^{-2\pi ka/b}\big)\Bigr),\quad k\in \mathbb{N}.$$



\smallskip



Уравнение \eqref{gim:eq-uravnmu1} в силу леммы \hyperlink{gim:l1}{1} имеет счетное множество действительных корней 

$\pm \mu_k$ и чисто мнимых корней $\pm i\mu_k,$ причем для положительных $\mu_k$ справедливо асимптотическое представление 

\begin{equation}

\mu_k=-\frac{\pi}{4}+\pi k+O(e^{-2\pi k}).

\label{gim:eq-muk}

\end{equation}



Тогда  $d$ может принимать значения

$d_k^{(1)}=\mu_k^2>0$ и $d_k^{(2)}=-\mu_k^2<0,$  и решениями задачи \eqref{gim:eq-uravn2}, \eqref{gim:eq-u6}, \eqref{gim:eq-uu6} будут соответственно функции

$$X_k^{(1)}(x)=\begin{cases}

\displaystyle{\frac{\sin[\mu_k (x-1)]}{\cos\mu_k }}, & x>0,\\[3mm] 

\displaystyle{\frac{\sh[\mu_k (x+1)]}{\ch\mu_k}},    & x<0,

\end{cases}\quad

X_k^{(2)}(x)=\begin{cases}

\displaystyle{\frac{\sh[\mu_k (x-1)]}{\ch\mu_k }},   & x>0,\\[3mm] 

\displaystyle{\frac{\sin[\mu_k (x+1)]}{\cos\mu_k}},  & x<0.

\end{cases}$$





При найденных $\mu_k$  решениями уравнения \eqref{gim:eq-uravn3} с учетом условий \eqref{gim:eq-u7} будут соответственно функции

\begin{equation}

Y_k^{(1)}(y)=

\begin{cases}

a_k^{(1)}\ch(y\sqrt{\mu_k^2-\lambda})+b_k^{(1)}\sh(y\sqrt{\mu_k^2-\lambda}), & y>0, \\

a_k^{(1)}\cos(y\sqrt{\mu_k^2-\lambda})+b_k^{(1)}\sin(y\sqrt{\mu_k^2-\lambda}), & y<0,

\end{cases} \label{gim:eq-ykk1}

\end{equation}

\begin{equation}

Y_k^{(2)}(y)=

\begin{cases}

a_k^{(2)}\cos(y\sqrt{\mu_k^2+\lambda})+b_k^{(2)}\sin(y\sqrt{\mu_k^2+\lambda}), & y>0, \\

a_k^{(2)}\ch(y\sqrt{\mu_k^2+\lambda})+b_k^{(2)}\sh(y\sqrt{\mu_k^2+\lambda}), & y<0,

\end{cases}

\label{gim:eq-ykk2}

\end{equation}

где $a_k^{(1)},$ $b_k^{(1)},$ $a_k^{(2)},$ $b_k^{(2)}$ "--- неизвестные пока коэффициенты.







Далее, удовлетворяя функции \eqref{gim:eq-ykk1} и \eqref{gim:eq-ykk2} граничным условиям \eqref{gim:eq-uu7}, получим системы для нахождения неизвестных коэффициентов $a_k^{(j)},$ $b_k^{(j)},$ $j=1, 2$:

\begin{equation}

\left\{\begin{array}{l}a_k^{(1)}\ch(\beta\sqrt{\mu_k^2-\lambda})+b_k^{(1)}\sh(\beta\sqrt{\mu_k^2-\lambda})=0, \\

a_k^{(1)}\cos(\alpha\sqrt{\mu_k^2-\lambda})-b_k^{(1)}\sin(\alpha\sqrt{\mu_k^2-\lambda})=0,\end{array}\right.\label{gim:eq-sistem1}

\end{equation}

\begin{equation}\left\{\begin{array}{l}a_k^{(2)}\cos(\beta\sqrt{\mu_k^2+\lambda})+b_k^{(2)}\sin(\beta\sqrt{\mu_k^2+\lambda})=0, \\

a_k^{(2)}\ch(\alpha\sqrt{\mu_k^2+\lambda})-b_k^{(2)}\sh(\alpha\sqrt{\mu_k^2+\lambda})=0.\end{array}\right.\label{gim:eq-sistem2}

\end{equation}



Системы \eqref{gim:eq-sistem1} и \eqref{gim:eq-sistem2} имеют нетривиальные решения тогда и только тогда, когда при всех $k\in \mathbb{N}$ соответственно определители этих систем равны 0:

\begin{multline}

\Delta_k^{(1)}(\alpha,\beta)=\cos(\alpha\sqrt{\mu_k^2-\lambda})  \sh(\beta\sqrt{\mu_k^2-\lambda})+

\\+

\sin(\alpha\sqrt{\mu_k^2-\lambda}) \ch(\beta\sqrt{\mu_k^2-\lambda})=0,\label{gim:eq-defk1}

\end{multline}

\begin{multline}

\Delta_k^{(2)}(\alpha,\beta)=\ch(\alpha\sqrt{\mu_k^2+\lambda}) \sin(\beta\sqrt{\mu_k^2+\lambda})+ \\ +

\sh(\alpha\sqrt{\mu_k^2+\lambda}) \cos(\beta\sqrt{\mu_k^2+\lambda})=0.\label{gim:eq-defk2}

\end{multline}



В этом случае 

$$b_k^{(1)}=a_k^{(1)}\ctg(\alpha\sqrt{\mu_k^2-\lambda}),\quad  

b_k^{(2)}=a_k^{(2)}\cth(\alpha\sqrt{\mu_k^2+\lambda}),$$ и функции \eqref{gim:eq-ykk1} и \eqref{gim:eq-ykk2} примут вид

$$

Y_k^{(1)}(y)= \begin{cases}

a_k^{(1)}\big(\cos(y\sqrt{\mu_k^2-\lambda})+\ctg(\alpha\sqrt{\mu_k^2-\lambda})\sin(y\sqrt{\mu_k^2-\lambda})\big), & y<0,\\

a_k^{(1)}\big(\ch(y\sqrt{\mu_k^2-\lambda})+\ctg(\alpha\sqrt{\mu_k^2-\lambda})\sh(y\sqrt{\mu_k^2-\lambda})\big), & y>0,

\end{cases}$$

$$

Y_k^{(2)}(y)= \begin{cases}

a_k^{(2)}\big(\ch(y\sqrt{\mu_k^2+\lambda})+\cth(\alpha\sqrt{\mu_k^2+\lambda})\sh(y\sqrt{\mu_k^2+\lambda})\big), &y<0,\\

a_k^{(2)}\big(\cos(y\sqrt{\mu_k^2+\lambda})+\th(\alpha\sqrt{\mu_k^2+\lambda})\sin(y\sqrt{\mu_k^2+\lambda})\big),& y>0,

\end{cases}$$

где $a_k^{(1)},$ $a_k^{(2)}$  "--- произвольные коэффициенты.



\vskip2mm



\hypertarget{gim:l2}{}



{\small \sc Лемма 2.} {\it Собственными значениями спектральной задачи $ D_\lambda$ являются числа

$$

\lambda_{k, 0}^{(1,1)}=\mu_k^2, \quad 

\lambda_{k, n}^{(1,2)}=\mu_k^2- (c_n^{(1)})^2,\quad 

\lambda_{k, n}^{(1,3)}=\mu_k^2+ (c_n^{(2)} )^2,$$

$$

\lambda_{k, 0}^{(2,1)}=-\mu_k^2,\quad 

\lambda_{k, n}^{(2,2)}=-\mu_k^2+ (c_n^{(2)} )^2,\quad 

\lambda_{k, n}^{(2,3)}=-\mu_k^2- (c_n^{(1)} )^2,$$

где $\mu_k$ "--- корни уравнения \eqref{gim:eq-uravnmu1}$,$ определяемые формулой} 

\eqref{gim:eq-muk}, $c_n^{(1)},$ $c_n^{(2)}$ "--- {\it корни уравнения $\tg(\alpha c)=-\th(\beta c),$ определяемые так}:

$$ c_n^{(1)}=-\frac{\pi}{4\alpha}+\frac{\pi}{\alpha}n+O\big(e^{-2\pi n \beta/\alpha}\big),

\quad 

c_n^{(2)}=-\frac{\pi}{4\beta}+\frac{\pi}{\beta}n+O\big(e^{-2\pi n \alpha/\beta}\big),\quad n=1, 2, \dots\,.$$



\smallskip

\proof

Собственные значения поставленной задачи являются корнями уравнений \eqref{gim:eq-defk1} и \eqref{gim:eq-defk2}.

Уравнение \eqref{gim:eq-defk1} равносильно уравнению

$$\tg(\alpha\sqrt{\mu_k^2-\lambda})=-\th(\beta\sqrt{\mu_k^2-\lambda}).$$

Обозначим $\sqrt{\mu^2_k-\lambda}=c,$ тогда уравнение

$\tg (\alpha c)=-\th (\beta c)$

согласно лемме \hyperlink{gim:l1}{1} имеет следующие корни:

$$

c_0=0,\qquad 

c_n^{(1)}=\Bigl(-\frac{\pi}{4\alpha}+\frac{\pi}{\alpha}n+O\big(e^{-2\pi n\beta/\alpha}\big)\Bigr),$$

$$ ic_n^{(2)}=i\Bigl(-\frac{\pi}{4\beta}+\frac{\pi}{\beta}n+O\big(e^{-2\pi n\alpha/\beta}\big)\Bigr),\quad 

c_n^{(3)}=-\Bigl(-\frac{\pi}{4\alpha}+\frac{\pi}{\alpha}n+O\big(e^{-2\pi n\beta/\alpha}\big)\Bigr),$$

$$ic_n^{(4)}=- i\Bigl(-\frac{\pi}{4\beta}+\frac{\pi}{\beta}n+O\big(e^{-2\pi n\alpha/\beta}\big)\Bigr),$$

где $n\in \mathbb{N}.$

Тогда собственными значениями будут действительные числа 

$$

\lambda_{k, 0}^{(1,1)}=\mu_k^2,\quad 

\lambda_{k, n}^{(1,2)}=\mu_k^2- (c_n^{(1)} )^2,\quad 

\lambda_{k ,n}^{(1,3)}=\mu_k^2+(c_n^{(2)})^2.$$

Аналогично из уравнения \eqref{gim:eq-defk2} найдем

$$

\lambda_{k, 0}^{(2,1)}=-\mu_k^2,\quad 

\lambda_{k, n}^{(2,2)}=-\mu_k^2+ (c_n^{(2)} )^2,\quad 

\lambda_{k, n}^{(2,3)}=-\mu_k^2- (c_n^{(1)} )^2.\quad \square$$









\vskip1mm



Соответствующие собственные функции имеют вид

$$

u^{(j,l)}_{k,n}(x,y)=X^{(j)}_k(x)Y^{(j)}_k(y,\lambda^{(j,l)}_{k,n}),$$

где

$$

Y^{(j)}_k(y,\lambda^{(j,l)}_{k,n})=Y^{(j)}_k(y)\big|_{\lambda=\lambda^{(j,l)}_{k,n}},\quad j=1, 2,\;\;l=1, 2, 3.

$$



Система $\bigl\{X_k^{(1)}(x),\;X_k^{(2)}(x)\bigr\}$  не ортогональна в $L_2[-1, 1].$

Задача, сопряженная  к задаче \eqref{gim:eq-uravn2}, \eqref{gim:eq-u6}, \eqref{gim:eq-uu6}, имеет следующий вид:

\begin{gather}

\mathop{\rm sgn} x \cdot  Z''+d \cdot  Z=0,\quad x\in (-1, 0)\cup (0, 1),\label{gim:eq-sopr1}

\\

Z(0-0)=-Z(0+0),\quad 

Z'(0-0)=-Z'(0+0),\quad 

 Z(-1)=Z(1)=0.\label{gim:eq-sopr3}

 \end{gather}

Решениями задачи \eqref{gim:eq-sopr1}, \eqref{gim:eq-sopr3} являются функции

$$

Z_k^{(1)}(x)=

\begin{cases}

-\displaystyle{\frac{\sin[\mu_k (x-1)]}{\cos\mu_k }}, &x>0,\\[3mm]

\displaystyle{\frac{\sh[\mu_k (x+1)]}{\ch\mu_k}}, & x<0,

\end{cases}\quad

Z_k^{(2)}(x)=

\begin{cases}

-\displaystyle{\frac{\sh[\mu_k (x-1)]}{\ch\mu_k }}, & x>0,\\[3mm]

\displaystyle{\frac{\sin[\mu_k (x+1)]}{\cos\mu_k}}, & x<0.

\end{cases}$$



Система $\bigl\{Z_k^{(1)};\, Z_k^{(2)}\bigr\}$ является биортогонально сопряженной с системой \linebreak $\bigl\{X_k^{(1)};\,X_k^{(2)}\bigr\},$ т.е. имеет место равенство

$$

\int _{-1}^1X_k^{(j)}(x)Z_m^{(l)}(x)dx=0,\quad k\ne m,\;j=1, 2;\;l=1, 2.

$$



Справедливо следующее утверждение.



\smallskip



\hypertarget{gim:l3}{}



{\small\sc Лемма 3.} {\it Система $\bigl\{Z_k^{(1)},\, Z_k^{(2)}\bigr\}$ полна в пространстве $L_2[-1, 1].$ }



\smallskip



{\bf 3. Единственность решения задачи $\boldsymbol D$.} 

Пусть теперь $\lambda\in \mathbb{R},$ $\lambda\ne\lambda_{k,n}^{i,j}.$



Далее рассмотрим функции

\begin{equation}

\begin{array}{r}

\displaystyle

u_k^{(1)}(y)=\int_{-1}^1u(x,y)Z_k^{(1)}(x)dx,\quad 

u_k^{(2)}(y)=\int_{-1}^1u(x,y)Z_k^{(2)}(x)dx,\\ 

k=1, 2, 3, \dots .

\end{array}

\label{gim:eq-uk}

\end{equation}



Вычислим вторую производную $\dfrac{d^2}{dy^2}\bigl(u_k^{(1)}(y)\bigr),$ преобразуем ее

на основании равенства \eqref{gim:eq-uravn1},

затем после двукратного интегрирования по частям получим равенство

$$

\mathop{\rm sgn}y\cdot (u_k^{(1)})''(y)+(\lambda-d) u_k^{(1)}(y)=0,

$$

совпадающее с уравнением \eqref{gim:eq-uravn3}. 

Таким образом, $u_k^{(1)}(y)\equiv Y_k^{(1)}(y),$ поэтому $u_k^{(1)}(y)$ определяются по формулам (\ref{gim:eq-ykk1}). Аналогичное справедливо и для функций $u_k^{(2)}(y),$ т.е. они определяются по формулам (\ref{gim:eq-ykk2}).



Тогда из граничных условий \eqref{gim:eq-u5} и равенств \eqref{gim:eq-uk} получим

\begin{equation}

\begin{cases}

\displaystyle

u_k^{(1)}(\beta)=

\int _{-1}^1 u(x,\beta)Z_k^{(1)}(x)dx=\int _{-1}^1 \varphi(x)Z_k^{(1)}(x)dx=\varphi_k^{(1)},\\

\displaystyle

u_k^{(2)}(\beta)=\int _{-1}^1\,u(x,\beta)Z_k^{(2)}(x)dx=

\int _{-1}^1\,\varphi(x)Z_k^{(2)}(x)dx=\varphi_k^{(2)},

\end{cases}

\label{gim:eq-uka}

\end{equation}

\begin{equation}

\begin{cases}

\displaystyle

u_k^{(1)}(-\alpha)= \int _{-1}^1 u(x,-\alpha)Z_k^{(1)}(x)dx=\int _{-1}^1 \psi(x)Z_k^{(1)}(x)dx=\psi_k^{(1)},\\

\displaystyle 

u_k^{(2)}(-\alpha)=\int _{-1}^1 u(x,-\alpha)Z_k^{(2)}(x)dx= \int _{-1}^1 \psi(x)Z_k^{(2)}(x)dx=\psi_k^{(2)}.

\end{cases}

\label{gim:eq-ukb}

\end{equation}



Тогда на основании \eqref{gim:eq-ykk1}, \eqref{gim:eq-ykk2}, \eqref{gim:eq-uka} и \eqref{gim:eq-ukb} получим систему для нахождения неизвестных коэффициентов $a_k^{(j)},$ $b_k^{(j)}$:

\begin{equation}\left\{\begin{array}{l}a_k^{(1)}\ch(\beta\sqrt{\mu_k^2-\lambda})+b_k^{(1)}\sh(\beta\sqrt{\mu_k^2-\lambda})=\varphi_k^{(1)}, \\

a_k^{(1)}\cos(\alpha\sqrt{\mu_k^2-\lambda})-b_k^{(1)}\sin(\alpha\sqrt{\mu_k^2-\lambda})=\psi_k^{(1)},\end{array}\right.\label{gim:eq-sistem11}

\end{equation}

\begin{equation}\left\{\begin{array}{l}a_k^{(2)}\cos(\beta\sqrt{\mu_k^2+\lambda})+b_k^{(2)}\sin(\beta\sqrt{\mu_k^2+\lambda})=\varphi_k^{(2)}, \\

a_k^{(2)}\ch(\alpha\sqrt{\mu_k^2+\lambda})-b_k^{(2)}\sh(\alpha\sqrt{\mu_k^2+\lambda})=\psi_k^{(2)}.\end{array}\right.\label{gim:eq-sistem12}

\end{equation}



Если при всех $k\in \mathbb{N}$ определители систем \eqref{gim:eq-sistem11} и \eqref{gim:eq-sistem12}, определяемые соотношениями  \eqref{gim:eq-defk1} и \eqref{gim:eq-defk2}, отличны от нуля,

то эти системы однозначно разрешимы:

$$a_k^{(1)}=\frac{1}{\Delta_k^{(1)}(\alpha,\beta)}\Big[\varphi_k^{(1)}\sin(\alpha\sqrt{\mu_k^2-\lambda})+\psi_k^{(1)}\sh(\beta\sqrt{\mu_k^2-\lambda})\Big],$$

$$b_k^{(1)}=\frac{1}{\Delta_k^{(1)}(\alpha,\beta)}\Big[\varphi_k^{(1)}\cos(\alpha\sqrt{\mu_k^2-\lambda})-\psi_k^{(1)}\ch(\beta\sqrt{\mu_k^2-\lambda})\Big],$$

$$a_k^{(2)}=\frac{1}{\Delta_k^{(2)}(\alpha,\beta)}\Big[\varphi_k^{(2)}\sh(\alpha\sqrt{\mu_k^2+\lambda})+\psi_k^{(2)}\sin(\beta\sqrt{\mu_k^2+\lambda})\Big],$$

$$b_k^{(2)}=\frac{1}{\Delta_k^{(2)}(\alpha,\beta)}\Big[\varphi_k^{(2)}\ch(\alpha\sqrt{\mu_k^2+\lambda})-\psi_k^{(2)}\cos(\beta\sqrt{\mu_k^2+\lambda})\Big].$$



Тогда с учетом найденных значений $a_k^{(j)},$ $b_k^{(j)}$ функции $u_k^{(j)}(y)$ примут вид

\begin{equation}

u_k^{(1)}(y)=

\left\{

\begin{array}{ll}\displaystyle 

\frac{1}{\Delta_k^{(1)}(\alpha,\beta)}\Big[\varphi_k^{(1)}\cdot

\Delta_k^{(1)}(\alpha,y)+ &

\\

\hspace{3cm}

+

\psi_k^{(1)}\cdot\sh\bigl((\beta-y)\sqrt{\mu_k^2-\lambda}\bigr)\Big], & y>0,\\

\displaystyle \frac{1}{\Delta_k^{(1)}(\alpha,\beta)}\Big[\varphi_k^{(1)}\cdot\sin\bigl((\alpha+y)\sqrt{\mu_k^2-\lambda}\bigr)+ &

\\

\hspace{3cm}

+

\psi_k^{(1)}\cdot\Delta_k^{(1)}(-y,\beta)\Big], & y<0,

\end{array}\right.\label{gim:eq-resh1}

\end{equation}

\begin{equation}

u_k^{(2)}(y)=\left\{\begin{array}{ll}

\displaystyle \frac{1}{\Delta_k^{(2)}(\alpha,\beta)} \Big[\varphi_k^{(2)}\cdot

\Delta_k^{(2)}(\alpha,y)+  &

\\

\hspace{3cm}

+

\psi_k^{(2)}\cdot\sin\bigl((\beta-y)\sqrt{\mu_k^2+\lambda}\bigr)\Big], & y>0,\\

\displaystyle \frac{1}{\Delta_k^{(2)}(\alpha,\beta)} \Big[\varphi_k^{(2)}\cdot\sh\bigl((\alpha+y)\sqrt{\mu_k^2+\lambda}\bigr)+ &

\\

\hspace{3cm}

+

\psi_k^{(2)}\cdot\Delta_k^{(2)}(-y,\beta)\Big], &  y<0,\end{array}\right.\label{gim:eq-resh2}

\end{equation}

где

$$\Delta_k^{(1)}(\alpha,y)=\cos(\alpha\sqrt{\mu_k^2-\lambda}) \sh(y\sqrt{\mu_k^2-\lambda})+

\sin(\alpha\sqrt{\mu_k^2-\lambda}) \ch(y\sqrt{\mu_k^2-\lambda}),$$

$$\Delta_k^{(1)}(-y,\beta)=\cos(y\sqrt{\mu_k^2-\lambda}) \sh(\beta\sqrt{\mu_k^2-\lambda})-\sin(y\sqrt{\mu_k^2-\lambda}) \ch(\beta\sqrt{\mu_k^2-\lambda}),$$

$$\Delta_k^{(2)}(\alpha,y)=\ch(\alpha\sqrt{\mu_k^2+\lambda}) \sin(y\sqrt{\mu_k^2+\lambda})+

\sh(\alpha\sqrt{\mu_k^2+\lambda}) \cos(y\sqrt{\mu_k^2+\lambda}),$$

$$\Delta_k^{(2)}(-y,\beta)=\ch(y\sqrt{\mu_k^2+\lambda}) \sin(\beta\sqrt{\mu_k^2+\lambda})-\sh(y\sqrt{\mu_k^2+\lambda}) \cos(\beta\sqrt{\mu_k^2+\lambda}).$$



Пусть $\varphi(x)=\psi(x)\equiv 0$ на $[-1, 1],$ тогда на основании \eqref{gim:eq-uka}, \eqref{gim:eq-ukb}, \eqref{gim:eq-resh1} и \eqref{gim:eq-resh2} получим

$$

\int _{-1}^1\, u(x,y)Z_k^{(j)}(x)dx=0,\quad j=1, 2, \; k=1, 2, 3, \dots.

$$





Отсюда в силу полноты системы $\bigl\{Z_k^{(1)}(x);\, Z_k^{(2)}(x)\bigr\}$ в $L_2[-1, 1]$ следует, что функция $u(x,y)=0$ при любом $y\in [-\alpha, \beta]$ почти всюду при $x\in[-1, 1].$ 

А в силу непрерывности $u(x,y)$ в $\overline D$ будет $u(x,y)\equiv 0$ в $\overline D.$



Пусть при некоторых $\alpha,$ $\beta$ и $k=p\in \mathbb{N}$ выполняется условие ${\Delta_{p}^{(1)}(\alpha, \beta)=0}$ или ${\Delta_{p}^{(2)}(\alpha,\beta)=0.}$ 

Пусть, например, $\Delta_{p}^{(1)}(\alpha,\beta)=0,$ а $\Delta_{p}^{(2)}(\alpha,\beta)\ne0.$ 

Тогда однородная задача \eqref{gim:eq-u2}--\eqref{gim:eq-u5}, где $\varphi(x)=\psi(x)=0$, 

имеет нетривиальное решение:

$$u_{p}(x,y)=\left\{\begin{array}{ll}

\displaystyle{\frac{\sin[\mu_{p}(x-1)]\Delta_{p}^{(1)}(\alpha,y)}{\cos\mu_{p}\cos(\alpha\sqrt{\mu^2_p-\lambda})}},

& (x,y)\in D_1,\\ 

\displaystyle{\frac{\sin[\mu_{p}(x-1)]\sin[(\alpha+y)\sqrt{\mu^2_p-\lambda}]}{\cos\mu_{p}\cos(\alpha\sqrt{\mu^2_p-\lambda})}}, & (x,y)\in D_2,\\ 

\displaystyle{\frac{\sh[\mu_{p}(x+1)]\sin[(\alpha+y)\sqrt{\mu^2_p-\lambda}]}{\ch\mu_{p}\cos(\alpha\sqrt{\mu^2_p-\lambda})}}, & (x,y)\in D_3,\\ 

\displaystyle{\frac{\sh[\mu_{p}(x+1)]\Delta_{p}^{(1)}(\alpha,y)}{\ch\mu_{p}\cos(\alpha\sqrt{\mu^2_p-\lambda})}}, & (x,y)\in D_4.\end{array}

\right.$$



Если при некотором $k=p\in \mathbb{N}$ имеем $\Delta_{k_0}^{(2)}(\alpha,\beta)=0,$ то также существует нетривиальное решение задачи.



Естественно возникает вопрос об обращении определителей $\Delta_{k}^{(j)}(\alpha,\beta)$ в нуль. Представим их в следующем виде:

\begin{equation}

\begin{array}{c}

\Delta_{k}^{(1)}(\alpha,\beta)=\Delta_k^{(1)}(\alpha,\beta)=\sqrt{\ch\big(2\beta\sqrt{\mu_k^2-\lambda}\big)}\sin\Big(\alpha\sqrt{\mu_k^2-\lambda}+\xi_k\Big),

\\

\Delta_{k}^{(2)}(\alpha,\beta)=\Delta_k^{(1)}(\alpha,\beta)=\sqrt{\ch\big(2\alpha\sqrt{\mu_k^2-\lambda}\big)}\sin\Big(\beta\sqrt{\mu_k^2-\lambda}+\chi_k\Big),

\end{array}

\label{gim:eq-d1}

\end{equation}

$$\xi_k=\arctg\left(\th\big(\beta\sqrt{\mu_k^2-\lambda}\big)\right),\quad  \chi_k=\arctg\left(\th\big(\alpha\sqrt{\mu_k^2-\lambda}\big)\right),$$

 причем 

 $$\lim\limits_{k\to\infty}\xi_k=\lim\limits_{k\to\infty}\chi_k=\pi/4,$$ и для всех $k\in\mathbb{N}$ справедливы неравенства $\xi_k<\pi/4,$ $\chi_k<\pi/4.$

Отсюда найдем  множества их нулей: 

\begin{equation}

\alpha_{k,m}=\frac{\pi m-\xi_k}{\sqrt{\mu_k^2-\lambda}},\quad 

\beta_{k,t}=\frac{\pi t-\chi_k}{\sqrt{\mu_k^2-\lambda}},\quad m,\;t,\;k\in \mathbb{N}.

\label{gim:eq-zn1}

\end{equation}



Равенства \eqref{gim:eq-zn1} представляют собой систему относительно $\alpha_{k,m}$ и $\beta_{k,t}.$ 

Можно убедиться, что эта система совместна, причем существует счетное множество её решений $(\alpha,\beta).$



Таким образом, справедливо следующее утверждение.



\smallskip



{\small\sc Теорема 1.} {\it Если существует решение задачи} \eqref{gim:eq-u2}--\eqref{gim:eq-u5}, {\it то оно единственно$,$ только если для всех $k\in \mathbb{N}$ выполняются условия $\Delta_{k}^{(j)}(\alpha,\beta)\ne 0,$ $j=1, 2.$ }



\vskip2mm



{\bf 4. Обоснование существования решения.}

Из формул \eqref{gim:eq-resh1} и \eqref{gim:eq-resh2} видно, что выражения $\Delta_{k}^{(j)}(\alpha,\beta)$ являются знаменателями дробей и при значениях $\alpha,$ $\beta,$ удовлетворяющих \eqref{gim:eq-zn1}, могут обратиться в нуль, т.е. возникает проблема <<малых знаменателей>>. 

Поэтому для обоснования существования решения задачи \eqref{gim:eq-u2}--\eqref{gim:eq-u5} необходимо показать существование чисел $\alpha$ и $\beta$ таких, что при больших $k$ выражения $\Delta_{k}^{(j)}(\alpha,\beta)$ отделены от нуля.



\smallskip





\hypertarget{gim:l4}{}



{\small \sc Лемма 4.} {\it Если выполнено одно из следующих условий}\/:

\begin{itemize}

\item[1)] {\it $\alpha$ "--- любое натуральное число$,$ кроме чисел вида $4p-3,$ $p\in \mathbb{N},\;$} 

\item[2)] {\it $\alpha$ "--- любое дробное число$,$ т.е. $\alpha=p/q,$ где $p,$ $ q\in \mathbb{N},$ $(p,q)=1,$ число $q-p$ не кратно $4,$}

\end{itemize} 

{\it   то существуют постоянная $ C_{01}>0$ и номер $k_{01}\in \mathbb{N}$ такие$,$ что для всех $k>k_{01}$ справедлива оценка }

\begin{equation}

|\Delta_{k}^{(1)}(\alpha,\beta)|\ge C_{01} e^{\pi k\beta}.

\label{gim:eq-ocenka1}

\end{equation}



\smallskip



\proof Оценим величину из \eqref{gim:eq-d1}: 

$$

\sqrt{\ch(2\beta\sqrt{\mu_k^2-\lambda})}\ge C  e^{\beta\sqrt{\mu_k^2-\lambda}}\ge \tilde C e^{\pi k\beta},

$$ где $\tilde C$ "--- некоторая положительная постоянная.



Так как

\begin{multline*}

\Bigl|\sin\Big(\alpha\sqrt{\mu_k^2-\lambda}+\xi_k\Big)-\sin\Bigl(\alpha(-\frac\pi 4+\pi k)+\frac\pi 4\Bigr)\Bigr|\le \\

\le \Bigl|\alpha\Bigl(\sqrt{\mu_k^2-\lambda}+\frac\pi 4-\pi k\Bigr)+\xi_k-\frac\pi 4\Bigr|\le 

\Bigl|\alpha(\pi k-\frac\pi 4-\sqrt{\mu_k^2-\lambda})\Bigr|=

\\

=\alpha\Bigl|

\frac{(\pi k-\pi/4)^2-\mu_k^2+\lambda}{\pi k-\pi/4+\sqrt{\mu_k^2-\lambda}}\Bigr|\le M_{1\alpha}k\,e^{-2\pi k},\quad M_{1\alpha}>0,

\end{multline*}

можем оценить

выражение 

$$

\Bigl|\sin\Bigl[\alpha\Bigl(-\frac{\pi}{4}+\pi k\Bigr)+\frac{\pi}{4}\Bigr]\Bigr|=|z_k|.

$$



\begin{itemize}

\item[1)] Пусть $\alpha\in \mathbb{N}.$ Возможны  три случая: $\alpha=2p,$ $\alpha=4p-1,$ $\alpha=4p-3,$ $p\in \mathbb{N}.$ 

В~первых двух случаях имеем $|z_k|={\rm const}> 0,$  а в третьем случае $|z_k|=0.$



\item[2)] Пусть $\alpha=p/q,$ $p, q\in \mathbb{N},$ $(p, q)=1.$  

Тогда 

$$z_k= \sin\Bigl[\frac{\pi}{4}\Bigl(1-\frac pq\Bigr)+\pi \frac{kp}{q}\Bigr].$$ Разделим число $kp$ на $q$ с остатком: $kp=sq+r,$ где $s, r\in \mathbb{N}\cup\{0\},$ $0\le r<q.$

Тогда

$$ z_k=\sin\frac{\pi(q-p+4r)}{4q}$$ и при условии, что $q-p$ не кратно 4, имеем $|z_k|={\rm const}> 0.$ $\square$

\end{itemize}





\smallskip



Нетрудно показать, что если $\alpha=4p-3,$ $p\in \mathbb{N},$ $p\ge 2,$ 

то существуют  постоянная $C'_{01}>0$ и номер $k'_{01}\in \mathbb{N}$ такие, что для всех $k>k'_{01}$  справедливы оценки

$$

\begin{array}{ll}

|\Delta_{k}^{(2)}(\alpha,\beta)|\le C'_{01} e^{-\pi k\beta} & \text{при}\;\; 0<\beta\le 1,\\

|\Delta_{k}^{(2)}(\alpha,\beta)|\le C'_{01} e^{\pi k(\beta-2)} & \text{при}\;\;\beta>1.

\end{array}

$$









Следовательно, приведенное в лемме \hyperlink{gim:l4}{4} условие $\alpha\ne 4p-3$ является существенным. 

Аналогично можно показать существенность условия, что  $q-p$ не кратно 4.



\smallskip



\hypertarget{gim:l5}{}



{\small\sc Лемма 5.} {\it Если выполнено одно из следующих условий}\/:

\begin{itemize}

\item[1)] {\it $\beta$ "--- любое натуральное число$,$ кроме чисел вида} $4p-3,$ $p\in \mathbb{N},$ 

\item[2)] $\beta$ "--- {\it любое дробное число$,$ т.е. $\beta=p/q,$ где $p, q\in \mathbb{N},$ $(p,q)=1,$ число $q-p$ не кратно} 4,

\end{itemize} 

{\it   то существуют постоянная $C_{02}>0$ и номер $ k_{02}\in \mathbb{N}$ такие$,$ что для всех $k>k_{02}$ справедлива оценка}

\begin{equation}

|\Delta_{k}^{(2)}(\alpha,\beta)|\ge C_{02} e^{\pi k\alpha}.

\label{gim:eq-ocenka2}

\end{equation}



\smallskip



Если выполнены оценки \eqref{gim:eq-ocenka1}, \eqref{gim:eq-ocenka2} и условия $\Delta_k^{(j)}(\alpha,\beta)\ne0$  при $k\le k_0=\max\{k_{01}, k_{02}\}$, то решение задачи \eqref{gim:eq-u2}--\eqref{gim:eq-u5} можно представить в виде суммы ряда Фурье

\begin{equation}

u(x,y)=\sum\limits_{k=1}^\infty u_k^{(1)}(y) X_k^{(1)}(x)+u_k^{(2)}(y) X_k^{(2)}(x).

\label{gim:eq-reshen}

\end{equation}



Определим условия на функции $\varphi(x)$ и $\psi(x),$ при которых ряд \eqref{gim:eq-reshen} будет сходиться равномерно в области $\overline D$ и допускать почленное дифференцирование по $x$ и $y.$



Рассмотрим следующие отношения:

$$

P_k^{(j)}(y)=\frac{\Delta_k^{(j)}(\alpha,y)}{\Delta_k^{(j)}(\alpha,\beta)},\quad

Q_k^{(1)}(y)=\frac{\sh\big((\beta-y)\sqrt{\mu_k^2-\lambda}\big)}{\Delta_k^{(1)}(\alpha,\beta)},$$

$$

Q_k^{(2)}(y)=\frac{\sin\big((\beta-y)\sqrt{\mu_k^2+\lambda}\big)}{\Delta_k^{(2)}(\alpha,\beta)}, \quad 

M_k^{(1)}(y)=\frac{\sin\big((\alpha+y)\sqrt{\mu_k^2-\lambda}\big)}{\Delta_k^{(1)}(\alpha,\beta)},

$$ 

$$

M_k^{(2)}(y)=\frac{\sh\big((\alpha+y)\sqrt{\mu_k^2+\lambda}\big)}{\Delta_k^{(2)}(\alpha,\beta)},\quad 

N_k^{(j)}(y)=\frac{\Delta_k^{(j)}(-y,\beta)}{\Delta_k^{(j)}(\alpha,\beta)},$$

где первые три выражения определены при $y>0,$ а последние три "--- при $y<0.$



\smallskip



{\small\sc Лемма 6.} {\it Пусть выполнены оценки} (\ref{gim:eq-ocenka1}), (\ref{gim:eq-ocenka2}) {\it при всех $k>k_0.$ Тогда для таких $k$ справедливы следующие оценки}\/:

$$|P_k^{(j)}(y)|\le C_1^{(j)},\quad |(P_k^{(j)}(y))'|\le C_2^{(j)} k,\quad |(P_k^{(j)}(y))''|\le C_3^{(j)} k^2,$$

$$|Q_k^{(1)}(y)|\le C_4^{(1)},\quad |(Q_k^{(1)}(y))'|\le C_5^{(1)} k,\quad |(Q_k^{(1)}(y))''|\le C_6^{(1)}  k^2,$$

$$|Q_k^{(2)}(y)|\le C_4^{(2)} e^{-\pi k\alpha},\;\; |(Q_k^{(2)}(y))'|\le C_5^{(2)} k e^{-\pi k\alpha},\;\; |(Q_k^{(2)}(y))''|\le C_6^{(2)} k^2 e^{-\pi k\alpha},$$

$$|M_k^{(1)}(y)|\le C_7^{(1)} e^{-\pi k\beta},\;\; |(M_k^{(1)}(y))'|\le C_8^{(1)}  k e^{-\pi k\beta},\;\; |(M_k^{(1)}(y))''|\le C_9^{(1)} k^2 e^{-\pi k\beta},$$

$$|M_k^{(2)}(y)|\le C_7^{(2)},\quad |(M_k^{(2)}(y))'|\le C_8^{(2)} k,\quad |(M_k^{(2)}(y))''|\le C_9^{(2)} k^2,$$

$$|N_k^{(j)}(y)|\le C_{10}^{(j)},\quad |(N_k^{(j)}(y))'|\le C_{11}^{(j)} k,\quad |(N_k^{(j)}(y))''|\le C_{12}^{(j)} k^2,$$

{\it где} $C_{l}^{(j)}>0,$ $j=1, 2,$ $l=\overline{1, 12}.$



\smallskip



\hypertarget{gim:l7}{}



{\small\sc Лемма 7.} {\it При всех $k>k_0$ справедливы оценки

$$|u_k^{(j)}(y)|\le C^{(j)}_{13}(|\varphi_k^{(j)}|+|\psi_k^{(j)}|),\quad |(u_k^{(j)}(y))'|\le C^{(j)}_{14}k(|\varphi_k^{(j)}|+|\psi_k^{(j)}|),$$

$$|(u_k^{(j)}(y))''|\le C^{(j)}_{15}k^2(|\varphi_k^{(j)}|+|\psi_k^{(j)}|),\quad y\in[-\alpha,\beta],$$

где}  $C^{(j)}_{l}>0,$ $j=1, 2,$ $l=\overline{13, 15}.$



\smallskip



В силу леммы \hyperlink{gim:l7}{7} ряд \eqref{gim:eq-reshen} и его первые производные в $\overline D,$ а вторые производные в $\overline D_i,$ $i=\overline {1, 4},$ по абсолютной величине мажорируются рядом

\begin{gather*}

C\sum\limits_{k=1}^\infty k^2 (|\varphi_k^{(1)}|+|\psi_k^{(1)}|+|\varphi_k^{(2)}|+|\psi_k^{(2)}|),\quad  

C={\rm const}>0.

\end{gather*}



\smallskip





\hypertarget{gim:l8}{}

{\small\sc Лемма 8.} {\it Если функции $\varphi(x),$ $\psi(x)\in C^1[-1, 1]\cap C^3[-1, 0]\cap C^3[0, 1]$ и на этом сегменте имеют кусочно-непрерывную производную четвертого порядка$,$ причем $$\varphi(-1)=\varphi(1)=\varphi''(-1)=\varphi''(1)=0,$$ 

$$\psi(-1)=\psi(1)=\psi''(-1)=\psi''(1)=0,$$ 

$$\varphi''(0+0)=-\varphi''(0-0),$$ 

$$\varphi'''(0+0)=-\varphi'''(0-0),$$  

$$\psi''(0+0)=-\psi''(0-0),$$ 

$$\psi'''(0+0)=-\psi'''(0-0),$$ то справедливы соотношения}

$$\varphi_k^{(1)}=\frac{p_k^{(1)}}{\mu_k^4},\qquad \varphi_k^{(2)}=\frac{p_k^{(2)}}{\mu_k^4},\qquad \psi_k^{(1)}=\frac{q_k^{(1)}}{\mu_k^4},\qquad

\psi_k^{(2)}=\frac{q_k^{(2)}}{\mu_k^4},$$

{\it где }

$$p_k^{(1)}=\int _{-1}^1 \varphi^{(4)}(x) Z_k^{(1)}(x)dx,\quad 

p_k^{(2)}=\int _{-1}^1 \varphi^{(4)}(x) Z_k^{(2)}(x)dx,$$ 

$$q_k^{(1)}=\int _{-1}^1 \psi^{(4)}(x) Z_k^{(1)}(x)dx,\quad 

q_k^{(2)}=\int _{-1}^1 \psi^{(4)}(x) Z_k^{(2)}(x)dx.$$



\smallskip



Для доказательства следует проинтегрировать по частям четыре раза  интегралы в равенствах (\ref{gim:eq-uka}) и (\ref{gim:eq-ukb}).



Можно показать, что для системы  $\bigl\{Z_k^{(1)}; Z_k^{(2)}\bigr\}$ справедливо неравенство Бесселя, и так как $\varphi^{(4)}(x)$ и $\psi^{(4)}(x)$ кусочно-непрерывны, то ряды $\sum_{k=1}^\infty (p_k^{(j)})^2$ и $\sum_{k=1}^\infty (q_k^{(j)})^2$ сходятся. 

Тогда сходится ряд 

$$\displaystyle{\sum\limits_{k=1}^\infty}\frac1k\bigl(|p_k^{(1)}|+|q_k^{(1)}|+|p_k^{(2)}|+|q_k^{(2)}|\bigr),$$

откуда следует равномерная сходимость ряда \eqref{gim:eq-reshen} и рядов, полученных дифференцированием по переменным $x$ и $y$ в $\overline D,$ а ряды из производных второго порядка сходятся в замкнутых областях $D_i,$ $i=\overline{1, 4}.$



Если при указанных в леммах \hyperlink{gim:l4}{4}  и 

\hyperlink{gim:l5}{5} числах $\alpha$ и $\beta$ и некоторых  

$k=k_1,$ $ k_2, \dots, k_l,$ где $1\le k_1<k_2<\cdots<k_l\le k_0,$ одно из выражений $\Delta_k^{(j)}(\alpha,\beta)=0$ (пусть для определенности $\Delta_k^{(1)}(\alpha,\beta)=0,$  $\Delta_k^{(2)}(\alpha,\beta)\ne 0$ при этих $k_i$),

то для разрешимости системы \eqref{gim:eq-sistem1} относительно $a_k^{(1)}$ и $b_k^{(1)}$ необходимо и достаточно выполнение условий

\begin{equation}

\varphi_{k_i}^{(1)}\cos (\alpha\sqrt{\mu_k^2-\lambda})=\psi_{k_i}^{(1)}\ch (\beta\sqrt{\mu_k^2-\lambda}), \quad i=\overline{1, l}.

\label{gim:eq-usl0}

\end{equation}





Тогда при $k=k_1, k_2, \dots, k_l$ получим

$$\tilde u_{k}(x,y)=\left\{\begin{array}{ll}

\displaystyle \frac{\sin[\mu_{k}(x-1)]}{\cos\mu_{k}}\Bigl(\ch(y\sqrt{\mu_k^2-\lambda})+ &

\\

\hspace{2cm}

\displaystyle

+

\frac{\varphi_k^{(1)}-

 \ch(\beta\sqrt{\mu_k^2-\lambda})}

{\sh(\beta\sqrt{\mu_k^2-\lambda})}\sh(y\sqrt{\mu_k^2-\lambda})\Bigr), & (x, y) \in D_1,

\\

\displaystyle

\frac{\sin[\mu_{k}(x-1)]} {\cos\mu_{k}}\Bigl(\cos(y\sqrt{\mu_k^2-\lambda})+ &

\\

\hspace{2cm}

\displaystyle

+

\frac{\varphi_k^{(1)}-\ch(\beta\sqrt{\mu_k^2-\lambda})}

{\sh(\beta\sqrt{\mu_k^2-\lambda})}\sin(y\sqrt{\mu_k^2-\lambda})\Bigr), & (x, y) \in D_2,

\\

\displaystyle\frac{\sh[\mu_{k}(x-1)]}{\ch\mu_{k}}\Bigl(\cos(y\sqrt{\mu_k^2-\lambda})+ 

\\

\hspace{2cm}

\displaystyle

+

\frac{\varphi_k^{(1)}-\ch(\beta\sqrt{\mu_k^2-\lambda})} {\sh(\beta\sqrt{\mu_k^2-\lambda})}\sin(y\sqrt{\mu_k^2-\lambda})\Bigr), & (x, y) \in D_3,\\

\displaystyle

\frac{\sh[\mu_{k}(x-1)]}{\ch\mu_{k}}\Bigl(\ch(y\sqrt{\mu_k^2-\lambda})+

\\

\hspace{2cm}

\displaystyle

+

\frac{\varphi_k^{(1)}-\ch(\beta\sqrt{\mu_k^2-\lambda})}

{\sh(\beta\sqrt{\mu_k^2-\lambda})}\sh(y\sqrt{\mu_k^2-\lambda})\Bigr), & (x, y) \in D_4.

\end{array}

\right.$$



Поэтому решение задачи \eqref{gim:eq-u2}--\eqref{gim:eq-u5} в этом случае определяется в виде суммы ряда

\begin{multline}

u(x,y)=\biggl(\sum\limits_{k=1}^{k_1-1}+\sum\limits_{k=k_1+1}^{k_2-1}+\cdots+\sum\limits_{k=k_l+1}^\infty\biggr) u_k^{(1)}(y) X_k^{(1)}(x)+

\\

+\sum\limits_{k=k_1,k_2,\dots,k_l}A_k \tilde u_k(x,y)+\sum\limits_{k=1}^\infty u_k^{(2)}(y) X_k^{(2)}(x),\label{gim:eq-reshen1}

\end{multline}

где $A_k$ "--- произвольные коэффициенты, причем конечные суммы следует считать равными нулю, если нижний предел суммирования больше верхнего.



Аналогичное решение строится в случае, когда  при некоторых $k$ будет $\Delta_k^{(2)}(\alpha,\beta)=0,$

 $\Delta_k^{(1)}(\alpha,\beta)\ne 0,$ или если оба знаменателя обращаются в нуль.



 \vskip1mm



Таким образом, доказано следующее утверждение.





\smallskip



{\small \sc Теорема 2.} 

{\it Пусть $\varphi(x),$ $\psi(x)$ удовлетворяют условиям леммы} \hyperlink{gim:l8}{8} {\it и выполнены оценки} \eqref{gim:eq-ocenka1}, \eqref{gim:eq-ocenka2}  {\it при $k>k_0$}. 

{\it Тогда если при указанных в леммах}~\hyperlink{gim:l4}{4}   {\it и} \hyperlink{gim:l5}{5}  {\it значениях $\alpha$ и $\beta$ при всех $k=\overline{1, k_0}$ выполнены условия} ${\Delta_k^{(1)}(\alpha,\beta)\ne0},$ $\Delta_k^{(2)}(\alpha,\beta)\ne0$, {\it то существует единственное решение задачи Дирихле} \eqref{gim:eq-u2}--\eqref{gim:eq-u5} {\it и оно определяется рядом} \eqref{gim:eq-reshen}.



{\it Если $\Delta_k^{(1)}(\alpha,\beta)=0$ при некоторых   $k=k_1,$  $k_2,\dots, k_l\le k_0,$ то задача} \eqref{gim:eq-u2}--\eqref{gim:eq-u5}  {\it разрешима только тогда$,$ когда выполнены условия } \eqref{gim:eq-usl0}, {\it и решение определяется в виде суммы ряда } \eqref{gim:eq-reshen1}.







\smallskip







\thanks{\textbf{Благодарности.}

Работа выполнена при поддержке Российского фонда фундаментальных исследований (проект № 14--01--97003-р\_Поволжье\_а).}





\thanks{{\bf ORCID}



{\bf Альфира Авкалевна Гималтдинова:}    \url{http://orcid.org/0000-0001-7535-213X}





}













\RusRef



\begin{thebibliography}{99}



\RBibitem{gimkur:add1-ru}

\by Гималтдинова~А.~А. 

\paper Задача Дирихле для уравнения смешанного типа с~двумя линиями перехода в прямоугольной области 

\pages 120--121

\inbook Четвертая международная конференция «Математическая физика и ее приложения»

\bookinfo материалы конф. 

\eds чл.-корр.~РАН И.~В.~Волович; д.ф.-м.н., проф. В.~П.~Радченко

\publaddr Самара

\publ СамГТУ

\yr 2014





\RBibitem{gim:bib:mon}

\by Сабитов~К.~Б., Биккулова~Г.~Г., Гималтдинова~А.~А.

\book К теории уравнений смешанного типа с~двумя линиями изменения типа

\publaddr Уфа

\publ Гилем

\yr 2006

\totalpages 150





\Bibitem{gim:bib:rasias}

\paper The Exterior Tricomi and Frankl Problems for Quaterelliptic-Quater\-hyper\-bolic Equations with Eight Parabolic Lines

\by Rassias~J.~M.

\jour Eur. J. Pure Appl. Math.

\yr 2011

\vol 4

\issue 2

\pages 186--208

\elink \url{http://www.ejpam.com/index.php/ejpam/article/view/1175/195}







\RBibitem{gim:bib:frankl}

\by Франкль Ф.~И.

\book Избранные труды по газовой динамике

\publaddr М.

\publ Наука

\yr 1973

\totalpages 711





\RBibitem{gim:bib:bits}

\paper Некорректность задачи Дирихле для уравнений смешанного типа

\by Бицадзе~А.~В.

\jour ДАН СССР

\yr 1958

\vol 122

\issue 2

\pages 167--170

\zmath{https://zbmath.org/?q=an:03234789}





\RBibitem{gim:bib:shabat}

\paper Примеры решения задачи Дирихле для уравнения смешанного типа

\by Шабат~Б.~В.

\jour ДАН СССР

\yr 1957

\vol 112

\issue 3

\pages 386--389

\zmath{https://zbmath.org/?q=an:0078.27803}





\Bibitem{gim:bib:can}

\paper A Dirichlet problem for an equation of mixed type with a discontinuous coefficient 

\by Cannon J.~R.

\jour  Annali di Matematica

\yr 1963

\vol 61

\issue 1 

\pages 371--377

\crossref{10.1007/bf02410656}





\RBibitem{gim:bib:nahushev}

\paper Критерий единственности задачи Дирихле для уравнения смешанного типа в цилиндрической области

\by  Нахушев А.~М.

\jour Диффер. уравн.

\yr 1970

\vol 6

\issue 1

\pages 190--191

\zmath{https://zbmath.org/?q=an:0214.10401}



\RBibitem{gim:bib:hach}

\paper О задаче Дирихле для одного уравнения смешанного типа

\by  Хачев~М.~М.

\jour Диффер. уравн.

\yr 1976

\vol 12

\issue 1

\pages 137--143





\RBibitem{gim:bib:sold1}

\paper Задачи типа Дирихле для уравнения Лаврентьева--Бицадзе. I: Теоремы единственности

\by  Солдатов~А.~П.

\jour Докл. РАН

\yr 1993

\vol 332

\issue 6

\pages 696--698

\mathnet{http://mi.mathnet.ru/dan5034}

\mathscinet{http://www.ams.org/mathscinet-getitem?mr=1272963}

%\transl

%\jour Dokl. Math.

%\yr 1994

%\vol 48

%\issue 2

%\pages 410--414



\RBibitem{gim:bib:sold2}

\paper Задачи типа Дирихле для уравнения Лаврентьева--Бицадзе. II: Теоремы существования

\by  Солдатов~А.~П.

\jour Докл. РАН

\yr 1993

\vol 333

\issue 1

\pages 16--18

\mathnet{http://mi.mathnet.ru/dan5027}

\mathscinet{http://www.ams.org/mathscinet-getitem?mr=1258440}

%\transl

%\jour Dokl. Math.

%\yr 1994

%\vol 48

%\issue 3

%\pages 433--437





\RBibitem{gim:bib:sab}

\by Сабитов~К.~Б.

\paper Задача Дирихле для уравнений смешанного типа в прямоугольной области

\jour Докл. РАН

\yr 2007

\vol 413

\issue 1

\pages 23--26

\mathnet{http://mi.mathnet.ru/dan701}

\mathscinet{http://www.ams.org/mathscinet-getitem?mr=2447063}







\RBibitem{gim:bib:sab-vag}

\paper Задача Дирихле для уравнения смешанного типа с двумя линиями вырождения в прямоугольной области

\by  Сабитов~К.~Б., Вагапова~Э.~В.

\jour Диффер. уравн.

\yr 2013

\vol 49

\issue 1

\pages 68--78





\RBibitem{gim:bib:iljin}

\by  Ильин~В.~А.

\paper Единственность и принадлежность $W^1_2$ классического решения смешанной задачи для самосопряженного гиперболического уравнения

\jour Матем. заметки

\yr 1975

\vol 17

\issue 1

\pages 91--101

\mathnet{http://mi.mathnet.ru/mz7227}

\mathscinet{http://www.ams.org/mathscinet-getitem?mr=369929}

\zmath{https://zbmath.org/?q=an:0315.35055}

%\transl

%\jour Math. Notes

%\yr 1975

%\vol 17

%\issue 1

%\pages 53--58

%\crossref{10.1007/BF01093843}





\RBibitem{gim:bib:arnold}

\by  Арнольд В.~И.

\paper Малые знаменатели.~I. Об отображениях окружности на себя

\jour Изв. АН СССР. Сер. матем.

\yr 1961

\vol 25

\issue 1

\pages 21--86

\mathnet{http://mi.mathnet.ru/izv3366}

\mathscinet{http://www.ams.org/mathscinet-getitem?mr=140699}

\zmath{https://zbmath.org/?q=an:0152.41905}

\morerref

\by Арнольд~В.~И.

\paper Исправления к~работе В.~Арнольда ``Малые знаменатели.~I''

\jour Изв. АН СССР. Сер. матем.

\yr 1964

\vol 28

\issue 2

\pages 479--480

\mathnet{http://mi.mathnet.ru/izv2961}

\mathscinet{http://www.ams.org/mathscinet-getitem?mr=164049}

\zmath{https://zbmath.org/?q=an:0135.42603}





\RBibitem{gim:bib:lomov2}

\paper Малые знаменатели в аналитической теории вырождающихся дифференциальных уравнений

\by  Ломов~И.~С.

\jour Диффер. уравн.

\yr 1993

\vol 29

\issue 12

\pages 2079--2089

\zmath{https://zbmath.org/?q=an:00789701}







\end{thebibliography}





\RuDate



\newpage



\makeentitle











\thanks{\textbf{Acknowledgments.}

This work has been supported by the Russian Foundation for Basic Research (project no. 14–01–97003-r\_Povolzh'e\_a).

}





\thanks{{\bf ORCID}



{\bf Alfira A. Gimaltdinova:}    \url{http://orcid.org/0000-0001-7535-213X}





}













\EngRef



\begin{thebibliography}{99}



\Bibitem{gimkur:add1-eng}

\lang In Russian

\by Gimaltdinova~A.~A. 

\paper The Dirichlet problem for mixed type equation  with two lines  of degeneracy  in a rectangular area

\pages 120--121

\inbook  The 4nd International Conference ``Mathematical Physics and its Applications''

\bookinfo Book of Abstracts and Conference Materials

\eds I.~V.~Volovich; V.~P.~Radchenko

\publaddr Samara

\publ Samara State Technical Univ.

\yr 2014

\lang In Russian





\Bibitem{gim:bib:mon:eng}

\lang In Russian

\by Sabitov~K.~B., Bikkulova~G.~G., Gimaltdinova~A.~A.

\book K teorii uravnenii smeshannogo tipa s~dvumia liniiami izmeneniia tipa \rm [On the Theory of Mixed Type Equations with Two Lines of Change of Type]

\publaddr Ufa

\publ Gilem

\yr 2006

\totalpages 150







\Bibitem{gim:bib:rasias:eng}

\paper The Exterior Tricomi and Frankl Problems for Quaterelliptic-Quater\-hyper\-bolic Equations with Eight Parabolic Lines

\by Rassias~J.~M.

\jour Eur. J. Pure Appl. Math.

\yr 2011

\vol 4

\issue 2

\pages 186--208

\elink \url{http://www.ejpam.com/index.php/ejpam/article/view/1175/195}







\Bibitem{gim:bib:frankl:eng}

\lang In Russian

\by Frankl~F.~I.

\book Izbrannye trudy po gazovoi dinamike \rm [Selected Works in Gas Dynamics]

\publaddr Moscow

\publ Nauka

\yr 1973

\totalpages 711





\Bibitem{gim:bib:bits:eng}

\lang In Russian

\by Bitsadze~A.~V.

\paper Incorrectness of Dirichlet's problem for the mixed type of equations in mixed regions

\jour Dokl. Akad. Nauk SSSR 

\yr 1958

\vol 122

\issue 2

\pages 167--170





\Bibitem{gim:bib:shabat:en}

\lang In Russian

\by Shabat~B.~V.

\paper Examples of solving the Dirichlet problem for equations of mixed type

\jour Dokl. Akad. Nauk SSSR 

\yr 1957

\vol 112

\issue 3

\pages 386--389





\Bibitem{gim:bib:can:eng}

\paper A Dirichlet problem for an equation of mixed type with a discontinuous coefficient 

\by Cannon J.~R.

\jour  Annali di Matematica

\yr 1963

\vol 61

\issue 1 

\pages 371--377

\crossref{10.1007/bf02410656}





\Bibitem{gim:bib:nahushev:eng}

\lang In Russian

\paper A criterion for uniqueness of the Dirichlet problem for a mixed-type equation in the cylindrical domain

\by  Nakhushev A.~M.

\jour Differ. Uravn.

\yr 1970

\vol 6

\issue 1

\pages 190--191

\zmath{https://zbmath.org/?q=an:0214.10401}



\Bibitem{gim:bib:hach:eng}

\lang In Russian

\paper On the Dirichlet Problem for an Equation of Mixed Type

\by  Khachev~M.~M.

\jour Differ. Uravn.

\yr 1976

\vol 12

\issue 1

\pages 137--143





\Bibitem{gim:bib:sold1:eng}

\by Soldatov~A.~P.

\paper Dirichlet-type problems for the Lavrent'ev--Bitsadze equation. I. Uniqueness theorems

\mathscinet{http://www.ams.org/mathscinet-getitem?mr=1272963}

\jour Dokl. Math.

\yr 1994

\vol 48

\issue 2

\pages 410--414



\Bibitem{gim:bib:sold2:eng}

\by Soldatov~A.~N.

\paper Dirichlet-type problems for the Lavrent'ev--Bitsadze equation. II. Existence theorems

\mathscinet{http://www.ams.org/mathscinet-getitem?mr=1258440}

\jour Dokl. Math.

\yr 1994

\vol 48

\issue 3

\pages 433--437





\Bibitem{gim:bib:sab:eng}

\by Sabitov~K.~B.

\paper The Dirichlet problem for equations of mixed type in a rectangular domain

\jour Dokl. Math. 

\vol 75 

\yr 2007

\issue  2

\pages 193–196 

\crossref{10.1134/S1064562407020056}

\mathscinet{http://www.ams.org/mathscinet-getitem?mr=2447063}







\Bibitem{gim:bib:sab-vag:eng}

\zmath{https://zbmath.org/?q=an:1268.35093}

\by Sabitov~K.~B., Vagapova~E.~V.

\paper Dirichlet problem for an equation of mixed type with two degeneration lines in a rectangular domain

\jour Differ. Equ.

\vol  49

\issue 1

\pages 68--78 

\yr 2013

\crossref{10.1134/s0012266113010072}



\Bibitem{gim:bib:iljin:eng}

\paper Proof of uniqueness and membership in $W_ 2^ 1$ of the classical solution of a mixed problem for a self-adjoint hyperbolic equation

\by Il'in~V.~A.

\jour Math. Notes

\yr 1975

\vol 17

\issue 1

\pages 53--58

\crossref{10.1007/BF01093843}





\Bibitem{gim:bib:arnold:eng}

\lang In Russian

\by Arnol'd~V.~I.

\paper Small denominators. I. Mapping the circle onto itself

\jour Izv. Akad. Nauk SSSR Ser. Mat.

\yr 1961

\vol 25

\issue 1

\pages 21--86

\mathscinet{http://www.ams.org/mathscinet-getitem?mr=140699}

\zmath{https://zbmath.org/?q=an:0152.41905}

\moreref

\lang In Russian

\by Arnol'd~V.

\paper Correction to V. Arnol'd's paper: ``Small denominators. I.''

\jour Izv. Akad. Nauk SSSR Ser. Mat. 

\yr 1964

\vol 28

\issue 2

\pages 479--480

\mathscinet{http://www.ams.org/mathscinet-getitem?mr=164049}

\zmath{https://zbmath.org/?q=an:0135.42603}





\Bibitem{gim:bib:lomov2:e}

\by Lomov~I.~S.

\paper Small denominators in analytic theory of degenerate differential equations

\jour Differ. Equ.

\vol 29

\issue 12

\pages 1811--1820 

\yr 1993

\zmath{https://zbmath.org/?q=an:00789701}







\end{thebibliography}



\EngDate





\newpage

\Russian



%\end{document}

